\documentclass{article}
\usepackage{amsmath, amssymb, amsthm}
\usepackage{hyperref}
\usepackage{geometry}
\geometry{margin=1in}

\title{Erd\H{o}s Problem \#1011}
\date{Last updated: 2025-09-10}
\author{Source: erdosproblems.com}

\begin{document}
\maketitle

\noindent\textbf{Status:} OPEN \quad \textbf{No prize}

\noindent\textbf{Tags:} graph theory

\medskip

\noindent\textbf{Problem Statement:}

Let $f_r(n)$ be minimal such that every graph on $n$ vertices with $\geq f_r(n)$ edges and chromatic number $\geq r$ contains a triangle. Determine $f_r(n)$.




Tur\'{a}n's theorem implies $f_2(n)=\lfloor n^2/4\rfloor+1$. Erd\H{o}s and Gallai \cite{Er62d} proved $f_3(n)=\lfloor \frac{1}{4}(n-1)^2\rfloor+2$. Simonovits showed in his PhD thesis (see the discussion on p. 358 of \cite{Si74}) that\[f_r(n)=\frac{n^2}{4}-\frac{g(r)}{2}{n}+O(1),\]where $g(r)$ is the largest $m$ such that, for any triangle-free graph with chromatic number $\geq r$, at least $m$ vertices of $G$ need to be removed to obtain a bipartite graph. Simonovits \cite{Si74} notes\[\frac{\log r}{\log\log r}r^2 \ll g(r) \ll (\log r)^2r^2.\]Hunter in the comments has noted that other results imply $g(r)\asymp r^2\log r$ - in fact\[(1/2-o(1))r^2\log r\leq g(r)\leq (2+o(1))r^2\log r.\]The lower bound follows from work of Davies and Illingworth \cite{DaIl22} (see [1104] ). The upper bound follows from work of Hefty, Horn, King, and Pfender \cite{HHKP25} on $R(3,k)$. Ren, Wang, Wang, and Yang \cite{RWWY24} showed that, for $n\geq 150$,\[f_4(n)=\left\lfloor\frac{(n-3)^2}{4}\right\rfloor+6.\]




References

[DaIl22] Davies, Ewan and Illingworth, Freddie, The {$\chi$}-{R}amsey problem for triangle-free graphs . SIAM J. Discrete Math. (2022), 1124--1134.

[Er62d] Erd\H{o}s, P., On a theorem of {R}ademacher-Tur\'{a}n . Illinois J. Math. (1962), 122--127.

[HHKP25] Z. Hefty, P. Horn, D. King, and F. Pfender, Improving $R(3,k)$ in just two bites . arXiv:2510.19718 (2025).

[RWWY24] S. Ren, J. Wang, S. Wang, and W. Yang, Extremal triangle-free graphs with chromatic number at least four . arXiv:2404.07486 (2024).

[Si74] Simonovits, M., Extermal graph problems with symmetrical extremal graphs. {A}dditional chromatic conditions . Discrete Math. (1974), 349--376.

\noindent\textbf{Related OEIS sequences:} possible


\bigskip
\noindent\small{Source: \url{https://www.erdosproblems.com/1011}}
\end{document}
